\begin{ctabular}{l|cccccc}
    \textbf{Hex} & A & B & C & D & E & F \\
    \textbf{Dec} & 10 & 11 & 12 & 13 & 14 & 15 \\
    \textbf{Bin} & 1010 & 1011 & 1100 & 1101 & 1110 & 1111
\end{ctabular}\vspace{-1em}

\section{Instruction Set Architecture (ISA)}
\label{subsec:isa_differences}

\begin{itemize}
    \item \textbf{ARM (32-bit) (Obsoleto):} 32 bit, massime prestazioni ma occupano maggior spazio in memoria.
    \item \textbf{Thumb (16-bit):} sottoinsieme 16 bit,
    Ideale per ridurre il footprint della memoria (circa 30\% di risparmio)
    \item \textbf{Thumb-2 (16/32-bit):} lunghezza variabile (16 e 32 bit). Permette di ottenere una densità di codice simile a Thumb 
    con prestazioni analoghe al set ARM a 32 bit.
    \item \textbf{RISC:} Istruzioni semplici e veloci eseguite in un ciclo di clock; usa un'architettura Load/Store (calcoli solo nei registri) e richiede più righe di codice Assembly.
    \item \textbf{CISC:} Istruzioni complesse e potenti che compiono più operazioni insieme; permette di eseguire calcoli direttamente sulla memoria RAM, riducendo le righe di codice Assembly.
\end{itemize}

\section{Logica della ALU e Program Status Register}
\begin{itemize}
    \item \textbf{N (Negative):} Se $MSB=1$, la flag è impostata a 1.
    \item \textbf{Z (Zero):} tutti i bit del risultato sono pari a zero.
    \item \textbf{C (Carry):} Indica un riporto oltre il bit più significativo. 
\begin{itemize}
    \item Fondamentale per operazioni \textit{unsigned}. \crd{Invertito per la sottrazione!}
    \item \textbf{Meccanismo \crd{sottrazione}:} $(A + \text{NOT}(B) + 1)$. \quad Zweierkompl.
    \begin{itemize}
        \item $C = 1 \implies A \ge B$ (Nessun prestito / Not-Borrow).
        \item $C = 0 \implies A < B$ (Prestito avvenuto / Borrow).
    \end{itemize}
\end{itemize}
    \item \textbf{V (oVerflow):} Indica un errore di segno nel calcolo in complemento 
    a due, Fondamentale per operazioni \textit{signed}.
    \begin{itemize}
        \item $(+) + (+) = (-)$ Impossibile
        \item $(-) + (-) = (+)$ Impossibile
    \end{itemize}
\end{itemize}

\subsubsection{Interpretazione del Segno e Salti Condizionati}
\label{subsec:sign_interpretation}

Poiché la ALU aggiorna simultaneamente sia la flag $C$ che la flag $V$,
la distinzione tra variabile \texttt{int} e \texttt{uint} non risiede nell'hardware, 
ma nel software.

